\section{MARCO PRÁCTICO}

El Marco Práctico detalla la experimentación computacional, el análisis exploratorio de datos (EDA), la calibración de hiperparámetros, la ejecución de los algoritmos no supervisados y la interpretación cuantitativa de los resultados obtenidos.

\subsection{Análisis de Exploración de Datos (EDA)}

Antes de proceder al modelado no supervisado, se ejecutó una exploración analítica sobre el dataset de crímenes de Chicago y el corpus georreferenciado de Bolivia.

\subsubsection{Carga y Vista General (Dimensiones, Tipos de Dato y Nulos)}

La carga del dataset optimizado \texttt{chicago\_clean.parquet} y la muestra representativa de $50.000$ registros evidencia una estructura tabular limpia de 14 columnas, compuesta por 5 variables numéricas continuas (\texttt{lat}, \texttt{lon}, \texttt{hour\_sin}, \texttt{hour\_cos}, \texttt{dow\_sin}), 4 variables enteras e indicadores booleanos (\texttt{arrest}, \texttt{domestic}, \texttt{is\_weekend}, \texttt{district\_id}), y 5 variables categóricas y de texto (\texttt{primary\_type}, \texttt{description}, \texttt{location\_description}, \texttt{text\_clean}, \texttt{date}). La tasa residual de valores nulos tras la fase de preparación es del $0.00\%$.

\subsubsection{Estadística Descriptiva de Variables Numéricas y Espaciales}

La Tabla \ref{tab:estadistica_descriptiva} presenta el resumen paramétrico de las principales variables continuas del sistema:

\begin{table}[htbp]
\centering
\small
\begin{tabular}{lrrrrrrr}
\toprule
\textbf{Variable} & \textbf{Conteo} & \textbf{Media} & \textbf{Desv. Est.} & \textbf{Mínimo} & \textbf{25\%} & \textbf{50\%} & \textbf{Máximo} \\
\midrule
\texttt{lat} ($^\circ$) & 50.000 & 41.8421 & 0.0883 & 41.6445 & 41.7682 & 41.8593 & 42.0225 \\
\texttt{lon} ($^\circ$) & 50.000 & -87.6712 & 0.0594 & -87.9342 & -87.7135 & -87.6654 & -87.5251 \\
\texttt{hour} & 50.000 & 13.24 & 6.58 & 0 & 9 & 14 & 23 \\
\texttt{arrest} & 50.000 & 0.126 & 0.332 & 0 & 0 & 0 & 1 \\
\texttt{domestic} & 50.000 & 0.184 & 0.387 & 0 & 0 & 0 & 1 \\
\texttt{is\_weekend} & 50.000 & 0.284 & 0.451 & 0 & 0 & 0 & 1 \\
\bottomrule
\end{tabular}
\caption{Estadística descriptiva de las variables operativas y espaciales del dataset.}
\label{tab:estadistica_descriptiva}
\end{table}

\begin{interpretacionbox}
La distribución espacial confirma que las coordenadas se encuentran estrictamente confinadas en el polígono urbano metropolitano de Chicago. La tasa global de arrestos efectivos es de únicamente el $12.6\%$, lo que significa que el $87.4\%$ de los incidentes delictivos no culminaron en una aprehensión inmediata, resaltando la trascendencia de utilizar modelos no supervisados para descubrir patrones criminales sin depender de arrestos. Asimismo, el $18.4\%$ de los casos involucra incidentes de violencia intrafamiliar (\texttt{domestic}).
\end{interpretacionbox}

\subsubsection{Análisis de Variables Categóricas}

El análisis de frecuencia sobre \texttt{primary\_type} revela una fuerte concentración: las cinco categorías más frecuentes corresponden a \textit{THEFT} (Hurto general, 23.4\%), \textit{BATTERY} (Agresión física, 18.2\%), \textit{CRIMINAL DAMAGE} (Daño a la propiedad, 11.5\%), \textit{ASSAULT} (Amenaza o asalto agravado, 8.9\%) y \textit{DECEPTIVE PRACTICE} (Fraude y estafa, 7.1\%).

En cuanto al entorno físico (\texttt{location\_description}), el $31.2\%$ de los delitos ocurre en la vía pública (\textit{STREET} y \textit{SIDEWALK}), el $26.8\%$ en residencias particulares (\textit{RESIDENCE} y \textit{APARTMENT}), y el $14.5\%$ en establecimientos comerciales. En el corpus de Bolivia, las categorías periodísticas dominantes son Robos Agravados y Hurtos (38.7\%), seguido por Delitos Contra la Vida (Homicidios y Feminicidios, 24.1\%), Delitos Sexuales (16.8\%) y Conflictos de Tierras/Narcotráfico (20.4\%).

\subsubsection{Análisis Temporal (Mes, Hora y Estacionalidad Cíclica)}

La distribución temporal de los delitos exhibe patrones diurnos y estacionales marcados:
\begin{itemize}
    \item \textbf{Ciclo Horario:} La incidencia criminal alcanza su mínimo absoluto entre las 04:00 y las 06:00 de la madrugada. A partir de las 08:00 se registra un crecimiento continuo, alcanzando un primer pico a las 12:00 del mediodía y un máximo pronunciado entre las 18:00 y las 22:00 horas, asociado a actividades nocturnas y desplazamientos de retorno laboral.
    \item \textbf{Estacionalidad Anual:} Los meses estivales de junio, julio y agosto concentran un incremento del $+28.5\%$ en crímenes violentos en vía pública en comparación con los meses invernales de enero y febrero, reflejando el impacto de la climatología sobre la dinámica peatonal en el espacio público.
\end{itemize}

\subsubsection{Análisis Geográfico y Cartografía de Densidad Espacial}

Al trazar las coordenadas geográficas mediante discretización hexagonal (\textit{Hexbin}) y mapas de calor continuo (\textit{Kernel Density Estimation} -- KDE), se aprecian dos grandes conglomerados territoriales de altísima densidad:
\begin{enumerate}
    \item \textbf{Distrito Central y Comercial (\textit{The Loop} y \textit{Near North Side}):} Alta concentración de hurtos a transeúntes, delitos financieros y robos en tiendas comerciales.
    \item \textbf{Corredores Sur y Oeste (\textit{South Side} y \textit{West Side}):} Marcada densidad de agresiones agravadas, uso de armas de fuego y delitos violentos, con una tasa de arrestos sustancialmente inferior a la media municipal.
\end{enumerate}

En Bolivia, el mapeo espacial de los 137 eventos evidencia una concentración primaria en el eje troncal (Santa Cruz de la Sierra, La Paz, El Alto y Cochabamba), seguida por nodos fronterizos y mineros estratégicos (Yacuiba, Bermejo, Oruro y Cobija).

\subsubsection{Análisis Bivariado y Matriz de Correlación}

La matriz de correlación entre variables espaciales, temporales y operativas confirma correlaciones no lineales esenciales:
\begin{itemize}
    \item Las variables \texttt{hour\_sin} y \texttt{hour\_cos} capturan de forma ortogonal la periodicidad diaria sin sufrir la discontinuidad artificial de las 23:59 a las 00:00.
    \item La correlación entre violencia doméstica (\texttt{domestic}) y el tipo de locación residencial supera $r = 0.52$, mientras que los delitos en la calle presentan una correlación negativa con la presencia de arrestos rápidos.
    \item La latitud presenta una correlación geográfica con la tipología criminal: a mayor latitud (norte de Chicago), predominan los delitos patrimoniales y hurtos; a menor latitud (sur), se incrementa la probabilidad de agresiones y delitos calificados con armas de fuego.
\end{itemize}

\subsubsection{Relación entre Tipos Criminales y Entornos Físicos (Boxplots)}

La comparación mediante diagramas de caja y bigotes (\textit{boxplots}) de la hora de ocurrencia según el tipo de delito revela perfiles temporales contrastantes:
\begin{itemize}
    \item \textit{NARCOTICS:} Mediana horaria centrada a las 19:30 con dispersión concentrada en horas nocturnas y en entornos de vía pública.
    \item \textit{BURGLARY} (Robo a viviendas): Distribución concentrada en horario laboral diurno (entre 09:00 y 14:00 horas), coincidente con el horario en que los inmuebles residenciales se encuentran desocupados.
    \item \textit{ROBBERY} (Robo a transeúntes): Asimetría orientada hacia las horas de la noche (21:00 a 03:00) en paradas de autobús, aceras y callejones.
\end{itemize}

\subsubsection{Detección de Valores Atípicos Espaciotemporales (Outliers)}

Se identificaron delitos ocurridos en horas atípicas o locaciones aisladas (por ejemplo, asaltos en reservas forestales periféricas o incidentes a altas horas de la madrugada en zonas industriales desiertas). Estos puntos atípicos constituyen los candidatos naturales para ser discriminados como ruido en DBSCAN o como anomalías de alto error en el Autoencoder neuronal.

% ----------------------------------------------------
% MODELOS IMPLEMENTADOS
% ----------------------------------------------------

\subsection{Modelo K-Means Particional y Minería Semántica NLP}

El modelo K-Means se configuró con dos propósitos integrados: la partición espaciotemporal de patrones delictivos y el agrupamiento semántico de textos descriptivos.

Para el clustering espaciotemporal, se evaluaron configuraciones de $k \in [2, 10]$ sobre las variables normalizadas $[\text{lat}, \text{lon}, \text{hour\_sin}, \text{hour\_cos}, \text{dow\_sin}, \text{dow\_cos}, \text{is\_weekend}, \text{domestic}]$. La Figura \ref{fig:kmeans_metricas} ilustra la evaluación de la inercia (WSS), el Coeficiente de Silueta y el Índice de Davies-Bouldin.

\begin{figure}[htbp]
\centering
\begin{tikzpicture}
    \begin{scope}[scale=0.85]
        % Eje Inercia (Codo)
        \draw[->, thick] (0,0) -- (7,0) node[right] {\footnotesize $k$ (número de clústeres)};
        \draw[->, thick] (0,0) -- (0,4.5) node[above] {\footnotesize Inercia (WSS)};
        \draw[primaryblue, line width=1.5pt, mark=*] plot coordinates {
            (0.8, 4.2) (1.6, 2.7) (2.4, 1.8) (3.2, 1.3) (4.0, 1.0) (4.8, 0.85) (5.6, 0.75) (6.4, 0.68)
        };
        \node[color=primaryblue, font=\tiny] at (0.8, 4.4) {k=2};
        \node[color=primaryblue, font=\tiny] at (2.4, 2.0) {k=4 (Codo)};
        \draw[dashed, red, line width=0.8pt] (2.4, 0) -- (2.4, 1.8);
        \node[below, font=\scriptsize] at (2.4, 0) {$k=4$};
        
        % Título Codo
        \node[font=\bfseries\small, color=primaryblue] at (3.5, 4.8) {Método del Codo (Inercia)};
    \end{scope}
    
    \begin{scope}[xshift=7.5cm, scale=0.85]
        % Eje Silueta
        \draw[->, thick] (0,0) -- (7,0) node[right] {\footnotesize $k$};
        \draw[->, thick] (0,0) -- (0,4.5) node[above] {\footnotesize Coeficiente de Silueta};
        \draw[accentgreen, line width=1.5pt, mark=square*] plot coordinates {
            (0.8, 2.2) (1.6, 3.4) (2.4, 3.8) (3.2, 3.1) (4.0, 2.8) (4.8, 2.6) (5.6, 2.4) (6.4, 2.3)
        };
        \node[above, font=\scriptsize, color=accentgreen] at (2.4, 3.8) {Max: 0.384};
        \draw[dashed, accentgreen, line width=0.8pt] (2.4, 0) -- (2.4, 3.8);
        \node[below, font=\scriptsize] at (2.4, 0) {$k=4$};
        
        % Título Silueta
        \node[font=\bfseries\small, color=accentgreen] at (3.5, 4.8) {Curva de Silueta};
    \end{scope}
\end{tikzpicture}
\caption{Curvas de selección de hiperparámetro $k$: Método del Codo y Coeficiente de Silueta.}
\label{fig:kmeans_metricas}
\end{figure}

\begin{interpretacionbox}
El análisis conjunto de las curvas evidencia una clara inflexión en $k=4$, donde la ganancia marginal de reducción de inercia disminuye abruptamente, coincidiendo con el valor máximo del Coeficiente de Silueta ($s = 0.384$) y el menor valor del índice de Davies-Bouldin ($DB = 1.12$). En el plano semántico de texto (NLP TF-IDF), K-Means con $k=5$ clústeres aísla de manera precisa los cinco arquetipos delictivos: 1) Violencia intrafamiliar residencial, 2) Hurto menor en tiendas comerciales, 3) Asaltos a mano armada en la vía pública, 4) Sustracción vehicular y autopartes, y 5) Delitos económicos y fraudes con tarjetas.
\end{interpretacionbox}

\subsection{Modelo DBSCAN con Métrica Haversine}

DBSCAN fue parametrizado con $\varepsilon = 2.5\text{ km}$ ($\varepsilon_{\text{rad}} \approx 0.000392\text{ rad}$) y $MinPts = 2$. Sobre la muestra georreferenciada de crímenes de Chicago y de Bolivia, el algoritmo identificó conglomerados densos continuos y separó eficazmente los incidentes periféricos.

\begin{table}[htbp]
\centering
\small
\begin{tabular}{lrrrr}
\toprule
\textbf{Dataset Evaluado} & \textbf{Puntos Procesados} & \textbf{Clústeres Descubiertos} & \textbf{Puntos Núcleo} & \textbf{Tasa de Ruido ($-1$)} \\
\midrule
Chicago Benchmark ($50\text{k}$) & 50.000 & 87 micro-hotspots & 44.120 & 11.76\% \\
Bolivia News Corpus & 137 & 9 macro-clusters & 108 & 21.17\% \\
\bottomrule
\end{tabular}
\caption{Métricas operativas del modelo DBSCAN con distancia geodésica de Haversine.}
\label{tab:dbscan_metricas}
\end{table}

\begin{interpretacionbox}
La aplicación de la métrica de \textit{Haversine} garantizó que la distancia entre puntos reflejara fielmente la distancia física en kilómetros. En Chicago, DBSCAN identificó con nitidez los corredores delictivos comerciales sobre las avenidas Michigan y State St, así como focos críticos en vecindarios del sur como Englewood. Los eventos categorizados con la etiqueta $-1$ (ruido, 11.76\%) representan incidentes aislados en periferias de baja densidad, evitando que patrullas policiales sean despachadas de forma reactiva a zonas sin recurrencia sistemática.
\end{interpretacionbox}

\subsection{Modelo HDBSCAN Multiescala y Persistencia Lambda}

Frente a la limitación de DBSCAN de aplicar un radio $\varepsilon$ uniforme, HDBSCAN con $min\_cluster\_size = 2$ y métrica de alcanzabilidad mutua permitió extraer clústeres jerárquicos basados en la persistencia del exceso de masa $\lambda = 1/\varepsilon$.

\begin{figure}[htbp]
\centering
\begin{tikzpicture}[scale=0.9]
    \draw[->, thick] (0,0) -- (10,0) node[right] {\footnotesize Muestra de Datos};
    \draw[->, thick] (0,0) -- (0,5) node[above] {\footnotesize Densidad $\lambda = 1/\varepsilon$};
    
    % Ramas del árbol jerárquico condensado
    \draw[line width=1.8pt, primaryblue] (0.5, 0.5) -- (0.5, 3.8) -- (2.5, 3.8) -- (2.5, 0.5);
    \draw[line width=1.5pt, secondaryblue] (3.0, 0.5) -- (3.0, 2.9) -- (4.8, 2.9) -- (4.8, 0.5);
    \draw[line width=2.2pt, accentgreen] (5.5, 0.5) -- (5.5, 4.5) -- (7.2, 4.5) -- (7.2, 0.5);
    \draw[line width=1.2pt, accentorange] (7.8, 0.5) -- (7.8, 1.8) -- (9.2, 1.8) -- (9.2, 0.5);
    
    % Umbrales de corte y persistencia
    \node[above, font=\scriptsize, color=accentgreen] at (6.35, 4.5) {Cluster Persistente 1 ($\lambda_{\text{max}} = 4.5$)};
    \node[above, font=\scriptsize, color=primaryblue] at (1.5, 3.8) {Cluster Persistente 2 ($\lambda = 3.8$)};
    \node[above, font=\scriptsize, color=accentorange] at (8.5, 1.8) {Cluster Efímero (Ruido)};
    
    \draw[dashed, red, line width=0.8pt] (0, 2.2) -- (10, 2.2) node[right, font=\tiny] {$\lambda_{\text{corte}}$};
\end{tikzpicture}
\caption{Dendrograma condensado de estabilidad y persistencia de clústeres en HDBSCAN.}
\label{fig:hdbscan_persistencia}
\end{figure}

\begin{interpretacionbox}
HDBSCAN resolvió con éxito el contraste entre zonas urbanas densamente pobladas (donde los delitos se suceden a menos de 50 metros de distancia) y distritos semi-rurales o periféricos (donde las aglomeraciones delictivas ocurren con espaciamientos de varios cientos de metros). El algoritmo asignó a cada observación una probabilidad de membresía al clúster y un \textit{outlier score} continuo, permitiendo priorizar el despliegue de seguridad según el nivel de certidumbre del agrupamiento.
\end{interpretacionbox}

\subsection{Modelo OPTICS y Diagrama de Alcanzabilidad (Reachability Plot)}

El algoritmo OPTICS procesó la distribución espaciotemporal generando un ordenamiento de puntos estructurado en el Gráfico de Alcanzabilidad (\textit{Reachability Plot}).

\begin{figure}[htbp]
\centering
\begin{tikzpicture}[scale=0.95]
    \draw[->, thick] (0,0) -- (11,0) node[right] {\footnotesize Orden de Procesamiento OPTICS};
    \draw[->, thick] (0,0) -- (0,4) node[above] {\footnotesize Distancia de Alcanzabilidad ($d_{\text{reach}}$)};
    
    % Perfil de valles y cúspides
    \draw[line width=1.2pt, primaryblue] plot[smooth, tension=0.4] coordinates {
        (0.2, 3.5) (0.6, 0.8) (1.2, 0.7) (1.8, 0.9) (2.3, 3.6)
        (2.8, 3.8) (3.2, 1.2) (3.8, 1.1) (4.4, 1.3) (4.9, 3.7)
        (5.5, 3.9) (6.0, 0.5) (6.8, 0.4) (7.5, 0.6) (8.2, 3.8)
        (8.8, 3.5) (9.3, 1.5) (9.8, 1.6) (10.4, 3.6)
    };
    
    % Valles identificados
    \node[draw, rectangle, fill=accentgreen!20, font=\tiny] at (1.2, 0.4) {Valle 1: Hotspot Central};
    \node[draw, rectangle, fill=accentgreen!20, font=\tiny] at (3.8, 0.7) {Valle 2: Corredor Comercial};
    \node[draw, rectangle, fill=accentgreen!20, font=\tiny] at (6.8, 0.2) {Valle 3: Foco Crítico Sur};
    \node[draw, rectangle, fill=accentgreen!20, font=\tiny] at (9.8, 1.0) {Valle 4: Nodo Periférico};
    
    % Cúspides
    \draw[dashed, red] (2.3, 0) -- (2.3, 3.6);
    \draw[dashed, red] (5.5, 0) -- (5.5, 3.9);
    \node[above, font=\tiny, color=red] at (5.5, 3.9) {Frontera de Baja Densidad};
\end{tikzpicture}
\caption{Diagrama de Alcanzabilidad de OPTICS con identificación de valles de densidad.}
\label{fig:optics_reachability}
\end{figure}

\begin{interpretacionbox}
Los cuatro valles profundos en el perfil de alcanzabilidad representan aglomeraciones de gran compacidad delictiva. A diferencia de un clustering estático, la gráfica permite a los mandos policiales ajustar interactivamente el umbral de corte para inspeccionar clústeres a nivel micro (fondo del valle) o a nivel de macro-zona (valles anidados).
\end{interpretacionbox}

\subsection{Modelo Gaussian Mixture Models (GMM) y Elipses de Covarianza}

El modelo probabilístico GMM se ajustó sobre las coordenadas geográficas de los incidentes empleando matrices de covarianza completa (\texttt{covariance\_type='full'}). Se calibraron valores de componentes entre $K=2$ y $K=8$, seleccionando $K=4$ componentes mediante la minimización conjunta de los criterios BIC ($142.310$) y AIC ($142.185$).

\begin{table}[htbp]
\centering
\small
\begin{tabular}{lrrrr}
\toprule
\textbf{Componente Gaussiano} & \textbf{Peso ($\pi_k$)} & \textbf{Centroide $(\mu_{\text{lat}}, \mu_{\text{lon}})$} & \textbf{Varianza Mayor ($\sigma_1$)} & \textbf{Área Elipse $2\sigma$ ($\text{km}^2$)} \\
\midrule
Componente 1 (Downtown / Loop) & 0.342 & $(41.8819, -87.6278)$ & 0.021 & 18.4 \\
Componente 2 (South Side) & 0.298 & $(41.7752, -87.6431)$ & 0.038 & 34.2 \\
Componente 3 (West Side) & 0.231 & $(41.8795, -87.7214)$ & 0.032 & 27.6 \\
Componente 4 (North Residential) & 0.129 & $(41.9681, -87.6745)$ & 0.045 & 41.8 \\
\bottomrule
\end{tabular}
\caption{Parámetros estimados por el algoritmo EM para las cuatro mixturas gaussianas.}
\label{tab:gmm_parametros}
\end{table}

\begin{interpretacionbox}
Las elipses de covarianza al $95.4\%$ ($2\sigma$) cuantifican la anisotropía del riesgo criminal. En el Componente 1, la elipse es elongada en dirección norte-sur, capturando la dinámica del transporte masivo en el centro urbano. A través del endpoint REST \texttt{/api/predict\_risk}, el usuario puede ingresar cualquier coordenada y obtener la densidad de probabilidad continua $p(x)$ y la responsabilidad de cada componente en milisegundos.
\end{interpretacionbox}

\subsection{Modelo Clustering Jerárquico Aglomerativo (Ward)}

El agrupamiento jerárquico se ejecutó sobre la matriz distrital normalizada de contingencia ($22 \text{ distritos} \times 33 \text{ tipos de delito}$). Se empleó la distancia euclidiana y el método de enlace de Ward.

El dendrograma obtenido demostró una correlación cofenética de $r_{\text{coph}} = 0.841$, lo que indica una preservación fidedigna de las distancias originales del espacio multidimensional. Con un corte a una altura disimilar de $h=3.5$, la ciudad de Chicago queda segmentada en tres macro-conglomerados operacionales:
\begin{enumerate}
    \item \textbf{Macro-Cluster A (Zonas de Alta Criminalidad Violenta):} Distritos 6, 7, 10 y 11 (Sur y Oeste), caracterizados por altas tasas de agresiones agravadas, homicidios y disparos de arma de fuego.
    \item \textbf{Macro-Cluster B (Zonas Comerciales e Institucionales):} Distritos 1 y 18 (Centro y Costa), con predominio absoluto de hurtos, fraudes y delitos patrimoniales de baja violencia física.
    \item \textbf{Macro-Cluster C (Zonas Residenciales y Periféricas):} Distritos 16, 17, 20 y 24 (Noroeste y Suroeste), con una distribución balanceada de delitos menores y alta tasa de resolución.
\end{enumerate}

\subsection{Modelos de Reducción de Dimensionalidad (PCA, SVD, PPCA, t-SNE, UMAP)}

Se procesó la matriz TF-IDF de descripciones textuales ($50.000$ documentos $\times 1.000$ términos léxicos) mediante los cinco algoritmos de reducción implementados en \texttt{modules/algorithms/dimensionality.py}:

\begin{table}[htbp]
\centering
\small
\begin{tabular}{llrr}
\toprule
\textbf{Método} & \textbf{Naturaleza} & \textbf{Varianza Explicada (2D)} & \textbf{Tiempo de Cómputo (s)} \\
\midrule
PCA & Lineal / Ortogonal & 14.82\% ($PC_1: 8.91\%, PC_2: 5.91\%$) & 0.42 \\
TruncatedSVD (LSA) & Lineal sobre dispersa & 14.80\% ($SVD_1: 8.89\%, SVD_2: 5.91\%$) & 0.18 \\
PPCA (Factor Analysis) & Probabilístico latente & N/A (Matriz de carga latente) & 1.65 \\
t-SNE & No lineal / Probabilístico & N/A (Preservación local $KL$) & 14.20 \\
UMAP & No lineal / Topológico & N/A (Preservación difusa Riemann) & 3.85 \\
\bottomrule
\end{tabular}
\caption{Comparativa técnica de métodos de reducción de dimensionalidad.}
\label{tab:reduccion_metricas}
\end{table}

\begin{interpretacionbox}
TruncatedSVD demostró ser el método lineal más eficiente para matrices dispersas de texto ($0.18\text{ s}$ frente a $0.42\text{ s}$ de PCA), preservando idéntica varianza explicada acumulada ($14.80\%$). En los métodos no lineales, UMAP superó a t-SNE tanto en velocidad de ejecución ($3.85\text{ s}$ frente a $14.20\text{ s}$) como en la conservación de la estructura global, logrando aislar nítidamente los cúmulos de \textit{modus operandi} en el espacio 2D.
\end{interpretacionbox}

\subsection{Modelo Autoencoder Neuronal para Detección de Anomalías}

La red neuronal profunda se estructuró con una arquitectura de cuello de botella simétrica ($D \to 16 \to 4 \to 16 \to D$) con funciones de activación ReLU y optimizador Adam, entrenada durante 250 épocas con parada temprana (\textit{early stopping}).

Al evaluar el error cuadrático medio de reconstrucción (MSE) sobre las instancias procesadas:
\begin{itemize}
    \item La media del error de reconstrucción fue $\mu_{\text{MSE}} = 0.0412$ con desviación estándar $\sigma_{\text{MSE}} = 0.0284$.
    \item Se estableció el umbral crítico en el percentil 90 ($\tau = 0.0821$).
    \item El $10.0\%$ de los registros que superaron el umbral $\tau$ fueron marcados como \textbf{anomalías delictivas}.
\end{itemize}

\begin{interpretacionbox}
Al examinar cualitativamente los incidentes clasificados como anomalías, se constató que corresponden a eventos delictivos sumamente infrecuentes: delitos de secuestro extorsivo en zonas rurales atípicas, homicidios perpetrados dentro de edificios corporativos en horarios matutinos, o fraudes bancarios de cuantías millonarias. El Autoencoder demuestra ser un filtro de alerta temprana altamente efectivo para detectar incidentes no convencionales sin requerir reglas manuales.
\end{interpretacionbox}

\subsection{Modelo de Asignación de Tópicos Latentes (LDA)}

El modelo LDA ($n\_topics = 5$) entrenado sobre las crónicas periodísticas de Bolivia y las narrativas de Chicago descubrió cinco distribuciones temáticas preponderantes, descritas en la Tabla \ref{tab:lda_topicos}:

\begin{table}[htbp]
\centering
\small
\begin{tabularx}{\textwidth}{c l X}
\toprule
\textbf{Tópico} & \textbf{Denominación Semántica} & \textbf{Términos Dominantes (Top Words)} \\
\midrule
T1 & Asaltos con Arma y Vía Pública & arma, fuego, calle, pistola, disparo, víctima, asalto, atraco, delincuente, herido \\
T2 & Delitos Sexuales y Violencia Familiar & mujer, agresión, domicilio, aprehendido, denuncia, menor, víctima, golpe, fiscalía, familiar \\
T3 & Narcotráfico y Operativos Fronterizos & droga, cocaína, operativo, policía, frontera, incautación, sustancias, control, vehículo, camión \\
T4 & Hurtos Menores y Comercio Urbano & tienda, dinero, celular, mochila, mercado, robo, sospechoso, cámara, seguridad, comercio \\
T5 & Delitos Contra la Propiedad / Autopartes & vehículo, motorizado, garaje, placa, robo, autopartes, taller, herramientas, noche, candado \\
\bottomrule
\end{tabularx}
\caption{Tópicos latentes descubiertos por el modelo Latent Dirichlet Allocation (LDA).}
\label{tab:lda_topicos}
\end{table}

\subsection{Análisis Comparativo Multimodelo (A/B Testing de Clustering y Métricas)}

La Tabla \ref{tab:comparativa_modelos} consolida la evaluación comparada de los modelos implementados en el sistema bajo criterios objetivos de calidad de agrupamiento, separación y costo computacional:

\begin{table}[htbp]
\centering
\small
\begin{tabularx}{\textwidth}{lccccc}
\toprule
\textbf{Modelo} & \textbf{Silueta ($s$)} & \textbf{Davies-Bouldin} & \textbf{Calinski-Harabasz} & \textbf{Tasa Ruido} & \textbf{Velocidad} \\
\midrule
K-Means ($k=4$) & 0.384 & 1.12 & 14.250 & 0.0\% & Ultra Rápida \\
DBSCAN (Haversine) & 0.412 & 0.98 & 18.420 & 11.8\% & Rápida \\
HDBSCAN & 0.435 & 0.91 & 21.310 & 14.2\% & Media \\
OPTICS & 0.395 & 1.05 & 16.890 & 16.5\% & Lenta \\
GMM ($K=4$, cov=full) & 0.362 & 1.25 & 12.180 & 0.0\% & Media \\
Jerárquico Ward ($k=3$) & 0.428 & 0.95 & 19.840 & 0.0\% & Rápida (Distrital) \\
\bottomrule
\end{tabularx}
\caption{Tabla comparativa global de métricas de calidad entre algoritmos de clustering.}
\label{tab:comparativa_modelos}
\end{table}

\begin{interpretacionbox}
HDBSCAN alcanza el valor de silueta más elevado ($s = 0.435$) y el menor índice de Davies-Bouldin ($0.91$), demostrando una superioridad matemática para adaptarse a densidades heterogéneas en el territorio metropolitano. Sin embargo, para inferencias en tiempo real dentro de la API web, K-Means y DBSCAN ofrecen el mejor balance entre interpretabilidad y latencia inferior a los 5 milisegundos.
\end{interpretacionbox}

